\section*{Alice and Bob vs Casino Challenge}

\subsection*{Objective}
Design a sequential circuit for Alice that responds to a betting game against a casino.

\subsection*{The Game}
Each round contains three bits:
\begin{itemize}
  \item $B$ for Bob,
  \item $C$ for the casino,
  \item $A$ for Alice, generated by your circuit.
\end{itemize}

Bob knows the full casino sequence ahead of time and may choose his own bit stream before the game begins. Alice cannot react verbally; her strategy must be embedded in the circuit itself.

\subsection*{Timing Model}
The validator runs exactly 7 clock cycles. The circuit is reset before each test, so all DFF state starts at 0.

On round $i$, the circuit receives the current Bob bit $B_i$ and casino bit $C_i$ and emits Alice's bit $A_i$.
Because the design is sequential, the information from round $i$ may also be stored and reused on round $i+1$.

\subsection*{Winning Condition}
A round is a win only when all three bits match:
\[
A_i = B_i = C_i
\]

The intended challenge is to build Alice so that, for every casino sequence, there exists a Bob strategy that wins at least 4 out of the 7 rounds.

\subsection*{Interface}
\begin{itemize}
  \item Inputs: $B$, $C$
  \item Output: $A$
\end{itemize}

\subsection*{Hints}
\begin{itemize}
  \item A single round is not enough; the circuit needs memory.
  \item A good design often stores a previous round's value in a DFF and uses it on the next cycle.
  \item If Bob can encode information about the casino sequence into his stream, Alice can try to decode it one cycle later.
\end{itemize}

\subsection*{Examples}
If Bob and the casino both send \texttt{1,1,1,1,1,1,1}, a sequential circuit with one-cycle memory can output
\[
0,1,1,1,1,1,1
\]
after the initial reset cycle.

\subsection*{Scoring Idea}
The puzzle is inspired by a minimax objective: for every casino sequence, search for the Bob sequence that gives Alice the best possible result.
The target is at least 4 wins out of 7.
